\begin{table}[!htbp]
\centering
\caption{Robustness to regional spillovers, HS3 LP: ATT by relative period}
\label{tab:ap_event_spillover_HS3_lp}
\vspace{-0.15cm}
\scriptsize
\setlength{\tabcolsep}{4pt}
\renewcommand{\arraystretch}{0.92}
\begin{tabular}{lccc}
\toprule
Relative period & Main & Exclude exposed DE & Exclude exposed district-DE \\
\midrule
-4 & 0.016 (0.046) & 0.053 (0.053) & 0.038 (0.053) \\
-3 & 0.014 (0.040) & -0.026 (0.040) & -0.009 (0.039) \\
-2 & 0.016 (0.038) & -0.003 (0.030) & -0.001 (0.029) \\
-1 & 0.000 & -0.014 (0.035) & -0.014 (0.038) \\
0 & 0.358 (0.045) & 0.350 (0.046) & 0.348 (0.045) \\
1 & 0.506 (0.046) & 0.485 (0.048) & 0.497 (0.046) \\
2 & 0.583 (0.047) & 0.586 (0.056) & 0.585 (0.047) \\
3 & 0.682 (0.048) & 0.658 (0.054) & 0.691 (0.049) \\
4 & 0.724 (0.058) & 0.736 (0.080) & 0.736 (0.059) \\
\bottomrule
\end{tabular}
\begin{flushleft}
\footnotesize Note: each cell reports the ATT and standard error in parentheses, in standard deviations of annual proficiency. Period \(k=-1\) is normalized as the baseline and therefore appears without a standard error.
\end{flushleft}
\end{table}



